\chapter{Human Capital and the Returns to Schooling}\label{ch:b3:human-capital}

People who stay longer in school earn more, in every country and every decade with
data. Whether the extra schooling causes the extra earnings is another question:
those who stay longer may be abler, better connected or more patient, and a diploma
may signal ability as much as it builds skill. This chapter treats schooling as an
investment, derives the earnings equation that measures its return, and follows
the natural experiments --- school-leaving laws, distance to a college, identical
twins --- that try to separate cause from selection.

\begin{openingchart}{Estimated return to a year of schooling, by least squares (OLS)
and by instrumental variables (IV) based on features of the school system, ten
studies (percent; IV with 95~\% intervals; the first model of each study).}{Card
(2000), NBER Working Paper 7769, Table 2}
\begin{tikzpicture}
\begin{axis}[omrecord, width=0.62\linewidth, height=0.5\linewidth, xmin=-2, xmax=28,
  ymin=0.3, ymax=10.7, xlabel={Percent per year of schooling},
  ytick={1,...,10}, yticklabels from table={figdata/bachelor-3/24-human-capital/out/estimates.csv}{study},
  y tick label style={font=\scriptsize}, legend pos=south east,
  legend style={font=\footnotesize}]
  \addplot[only marks, mark=square*, omThm, error bars/.cd, x dir=both, x explicit]
    table[x=iv, y=pos, x error expr=1.96*\thisrow{iv_se}] {figdata/bachelor-3/24-human-capital/out/estimates.csv};
  \addplot[only marks, mark=*, omDef] table[x=ols, y=pos] {figdata/bachelor-3/24-human-capital/out/estimates.csv};
  \legend{IV, OLS}
\end{axis}
\end{tikzpicture}
\end{openingchart}

% figdata: estimates_stats
In nine of the ten studies the IV estimate is above the OLS one; the median IV
estimate is 10.6~\% per year, the median OLS estimate 7.2~\%. If the OLS estimates were
inflated by ability, the IV estimates should be lower. They are not, and the
intervals are wide. Reading this pattern is the work of the chapter.

\begin{recall}
The grades 10--12 volume introduced \omterm{def:b3:human-capital:human-capital}{human capital}; the Year 2 volume built the
signalling model of diplomas, and multiple regression with omitted-variable bias.
\omterm{def:b3:instrumental-variables:instrument}{Instrumental variables} and the \omterm{def:b3:instrumental-variables:late}{local average treatment effect} are chapter 1's,
regression discontinuity chapter 4's.
\end{recall}

\begin{notation}[Symbols of this chapter]\label{not:b3:human-capital:symbols}
$S$ is years of schooling; $Y(S)$ annual earnings with $S$ years; $\rho$ the Mincer coefficient;
$A$ ability; $X$ years of experience.
\end{notation}

\section{The schooling decision}\label{sec:b3:human-capital:decision}

\begin{definition}[Human capital (investment model)]\label{def:b3:human-capital:human-capital}
\index{human capital}
\emph{Human capital} is the stock of skills that raises a person's productivity,
built by investment --- schooling, training --- whose cost is mostly the earnings
given up while learning, and whose return is higher earnings later.
\end{definition}

\begin{definition}[Internal rate of return to schooling]\label{def:b3:human-capital:irr}
\index{internal rate of return to schooling}
The \emph{internal rate of return} to a year of schooling is the discount rate at
which the present value of the extra earnings it brings equals the present value
of its costs.
\end{definition}

\begin{proposition}[The stopping rule and the Mincer coefficient]\label{prop:b3:human-capital:stopping}
Suppose schooling has no tuition, a person with $S$ years earns $Y(S)$ every year from $S$ for
ever, and she discounts at rate $r$. She chooses $S$ where $Y'(S)/Y(S) = r$. If people who choose
different $S$ are equally well off, then $\ln Y(S+1) - \ln Y(S) = r$: the coefficient $\rho$ of $\ln Y$ on $S$ is the
internal rate of return.
\end{proposition}

\begin{proof}
The present value is $\int_S^\infty Y(S)e^{-rt}\,dt = Y(S)e^{-rS}/r$. Its logarithm, $\ln Y(S) - rS - \ln r$, has derivative
$Y'(S)/Y(S) - r$, zero at the optimum. If $S$ and $S + 1$ give the same present value, $Y(S+1)e^{-r(S+1)} =
Y(S)e^{-rS}$, so $\ln Y(S+1) - \ln Y(S) = r$; the discount rate that equates the two streams --- the internal
rate of return --- is the log earnings gain.
\end{proof}

\begin{omfigure}
\begin{tikzpicture}[econ, x=0.5cm, y=0.5cm]
  \draw[econ axes] (0,0) -- (13,0) node[right, font=\scriptsize] {age};
  \draw[econ axes] (0,0) -- (0,7) node[above, font=\scriptsize] {earnings};
  \draw[curve] (2,3) -- (12.5,3) node[right, font=\scriptsize] {$Y(S)$};
  \draw[curve, omDef] (3,4.2) -- (12.5,4.2) node[right, font=\scriptsize] {$Y(S+1)$};
  \fill[omThm!35] (2,0) rectangle (3,3);
  \fill[omDef!15] (3,3) rectangle (12.5,4.2);
  \node[font=\scriptsize, align=center] at (2.5,-0.9) {cost};
  \node[font=\scriptsize] at (7.7,3.6) {gain every year};
  \draw[guide] (2,0) -- (2,3); \draw[guide] (3,0) -- (3,4.2);
\end{tikzpicture}
\par\smallskip{\small A year more of school: earnings $Y(S)$ given up for a year (the
cost), and $Y(S+1) - Y(S)$ gained every year after. The return is the rate that makes the
two worth the same today.}
\end{omfigure}

\section{The earnings equation}\label{sec:b3:human-capital:mincer}

\begin{definition}[Mincer earnings equation]\label{def:b3:human-capital:mincer}
\index{Mincer earnings equation}
The \emph{Mincer earnings equation} is the regression $\ln Y = a + \rho S + \beta_1X + \beta_2X^2 + u$ of log earnings on
years of schooling and a quadratic in experience.
\end{definition}

\begin{omchart}{Mean log annual earnings by years of schooling, men in the U.S.
Current Population Survey comparison sample, 1978, positive earnings: raw means
(grey) and means adjusted to the sample's mean experience with the fitted experience
terms (dots), which are comparable with the Mincer fit at mean experience (line).
Earnings are top-coded at 25{,}565 dollars, which binds for 30~\% of
the sample.}{Dehejia--Wahba / LaLonde data (cps\_controls); our regression}
\begin{tikzpicture}
\begin{axis}[omrecord, width=0.58\linewidth, height=0.38\linewidth, xmin=5.5,
  xmax=18.5, ymin=9.1, ymax=10.1, xlabel={Years of schooling},
  ylabel={Mean log earnings}]
  \addplot[only marks, mark=o, black!45, mark size=1.6pt] table[x=education, y=lw] {figdata/bachelor-3/24-human-capital/out/mincer_means.csv};
  \addplot[only marks, mark=*, omDef] table[x=education, y=adj] {figdata/bachelor-3/24-human-capital/out/mincer_means.csv};
  \addplot[omThm, thick, mark=none] table[x=education, y=lw] {figdata/bachelor-3/24-human-capital/out/mincer_fit.csv};
\end{axis}
\end{tikzpicture}
\end{omchart}

% figdata: mincer_stats
In 13{,}820 men with positive earnings, a year of schooling goes with log earnings
higher by 0.063 (s.e.\ 0.003) holding experience fixed, and earnings rise with experience
at a falling rate, peaking after about 27 years. The raw means fall from 6 to 10 years
of schooling (9.39 to 9.19) because the men with little schooling are older and more
experienced; held at mean experience, they rise. Two features of the data stand out: a
jump at 12 years, the high-school diploma, of 0.20 in the adjusted means against a
median step of 0.01 elsewhere; and the top-coding, which flattens the dots at the top,
below the straight line.

\begin{model}{The Mincer earnings equation}
\assumes Schooling is an investment with a constant return; earnings grow with
experience as workers keep investing in skills, at a falling rate.
\predicts Log earnings rise linearly in years of schooling and concavely in
experience.
\tested The linear form fails a direct test in the chart's sample: replacing the
single schooling term by one dummy per year of schooling improves the fit by far more
than chance ($F = 5.8$, $p < 0.001$), mostly through the jump at 12 years.
% figdata: mincer_stats
The size of the average return is well measured: across 1{,}120 estimates in 139
countries it is about 9~\% and very stable over decades, higher for women and in
low-income countries. % ledger: b3hc:pp
\breaks The years that complete a diploma pay more than the others (a sheepskin
effect, the Year 2 volume's signal); returns differ across people; and in many rich
economies the return to higher education has risen, making the relation convex.
\end{model}

\section{Ability bias and natural experiments}\label{sec:b3:human-capital:ability}

\begin{definition}[Ability bias]\label{def:b3:human-capital:ability-bias}
\index{ability bias}
\emph{Ability bias} is the omitted-variable bias in the OLS estimate of the return
to schooling that arises because ability raises earnings and is correlated with
schooling.
\end{definition}

\begin{proposition}[The sign of ability bias]\label{prop:b3:human-capital:bias}
If $\ln Y = a + \rho S + \gamma A + u$ with $u$ uncorrelated with $S$ and $A$, the OLS slope of $\ln Y$ on $S$ alone converges
to $\rho + \gamma\,\mathrm{cov}(S, A)/\mathrm{var}(S)$: upward if ability raises earnings ($\gamma > 0$) and the able study longer.
\end{proposition}

\begin{proof}
The OLS slope is $\mathrm{cov}(S, \ln Y)/\mathrm{var}(S) = [\rho\,\mathrm{var}(S) + \gamma\,\mathrm{cov}(S, A)]/\mathrm{var}(S)$, since $u$ is uncorrelated
with $S$.
\end{proof}

Two forces push the other way. Schooling is reported with error, which biases OLS
toward zero. And an instrument moves the schooling of a particular group --- those
who would have left at the minimum age, or who would not have gone to college had
it been far --- and if their return is higher than average, IV estimates their
\emph{local} return, the LATE of chapter 1.

\begin{record}[Compulsory schooling laws]
Britain raised its minimum school-leaving age in 1947 and 1973. Using both changes
as instruments for men's schooling, an early study found an IV return of 15.3~\% against
an OLS one of 6.1~\%. % ledger: b3hc:card.hw
A later study of the 1947 change also found about 15~\%; reanalysing the same survey gives
about 3~\% on average, none for women and 4--7~\% for men, and a panel with better
earnings data gives zero for women and 3--4~\% for men. % ledger: b3hc:dh
One law change, one survey and two teams gave answers five times apart.
\end{record}

\begin{measure}[Twins]
Identical twins share their genes and their upbringing, so a difference in their
schooling is not a difference in inherited ability. A survey of twins that asked
each twin about both their own and their twin's schooling found no upward \omterm{def:b3:human-capital:ability-bias}{ability
bias}, but a downward bias from reporting error; corrected for it, a year of
schooling raised wages by 16~\%. % ledger: b3hc:ak
Differencing removes the shared ability and doubles the weight of reporting
error, which is why the correction matters.
\end{measure}

\section{Signal or skill}\label{sec:b3:human-capital:signal}

The signalling model of the Year 2 volume predicts the same correlation between
schooling and earnings as the investment model. It differs in what the extra
earnings pay for: skill in one case, information about the worker in the other. Its
sharper predictions are the jump at a diploma, which the CPS data show at 12 years,
and returns that fall as employers learn workers' productivity on the job. For
policy the distinction matters through the social return: a year of school that
only sorts workers raises the private return without raising output.

\section{School resources}\label{sec:b3:human-capital:resources}

\begin{definition}[Education production function]\label{def:b3:human-capital:education-production}
\index{education production function}
An \emph{education production function} relates pupils' measured achievement to
school inputs --- class size, teachers, spending --- and to the pupils' own and family
characteristics.
\end{definition}

\begin{record}[Project STAR]
In Tennessee, 11{,}600 pupils and their teachers were randomly assigned from 1985--86
to small classes (13--17 pupils), regular classes (22--25) or regular classes with an
aide. Test scores rose by about 4 percentile points in a pupil's first year in a small
class and by about one point a year after; aides made little difference; the
effect was larger for minority pupils and those on free lunches. About 10~\% of pupils
moved between class types, and the analysis checks that this does not drive the
result. % ledger: b3hc:star
\end{record}

\section{Training on the job}\label{sec:b3:human-capital:training}

\begin{definition}[General and specific training]\label{def:b3:human-capital:training}
\index{general training}\index{specific training}
\emph{General training} raises a worker's productivity equally in many firms;
\emph{specific training} raises it only in the firm that provides it.
\end{definition}

\begin{proposition}[Who pays for training]\label{prop:b3:human-capital:becker}
In competitive labour markets, workers pay for \omterm{def:b3:human-capital:training}{general training}, through a lower
wage while trained; firms can recover, and so share, the cost of \omterm{def:b3:human-capital:training}{specific training}.
\end{proposition}

\begin{proof}
After \omterm{def:b3:human-capital:training}{general training} a worker's productivity rises by the same amount in every
firm, and competition among firms raises her wage by that amount: the training firm
gains nothing after training, so it provides training only if the wage during
training falls by its cost. After \omterm{def:b3:human-capital:training}{specific training} other firms' offers are
unchanged, so the training firm can pay less than the worker's new productivity
and recover its outlay; if it paid for all of it, the worker could quit at no loss,
so the two share the cost and the return.
\end{proof}

\section{Exercises}\label{sec:b3:human-capital:exercises}

\begin{exercise}[$\star$]\label{exo:b3:human-capital:1}
In the opening chart, which study has the largest gap between IV and OLS, and which
is the only one with IV below OLS?
\end{exercise}

\begin{exercise}[$\star$]\label{exo:b3:human-capital:2}
A one-year course costs 20{,}000 dalens of lost earnings and raises earnings by 1{,}500 dalens a
year for ever. Find its internal rate of return.
\end{exercise}

\begin{exercise}[$\star$]\label{exo:b3:human-capital:3}
A Mincer coefficient is 0.11. Convert it into a percentage earnings premium per year.
\end{exercise}

\begin{exercise}[$\star$]\label{exo:b3:human-capital:4}
What would twin data show if \omterm{def:b3:human-capital:ability-bias}{ability bias} inflated the OLS estimate?
\end{exercise}

\begin{exercise}[$\star\star$]\label{exo:b3:human-capital:5}
Estimate the Mincer equation in the CPS comparison sample, then add indicators for
black, Hispanic and married. How does the schooling coefficient change, and why?
\end{exercise}

\begin{exercise}[$\star\star$]\label{exo:b3:human-capital:6}
Each twin reports schooling with an independent error. Explain why the bias toward
zero is larger in the twin difference than in the OLS estimate across people.
\end{exercise}

\begin{exercise}[$\star\star$]\label{exo:b3:human-capital:7}
Cohorts born just after a school-leaving reform have 0.3 more years of schooling and
log earnings higher by 0.024 than those born just before. Find the Wald estimate.
\end{exercise}

\begin{exercise}[$\star\star$]\label{exo:b3:human-capital:8}
Who pays for a forklift licence, and who for training in a firm's own software? Use
the proposition.
\end{exercise}

\begin{exercise}[$\star\star$]\label{exo:b3:human-capital:9}
Under pure \omterm{def:b3:human-capital:human-capital}{human capital}, what should a regression of earnings on years of
schooling and a diploma indicator show?
\end{exercise}

\begin{exercise}[$\star\star\star$]\label{exo:b3:human-capital:10}
Prove the Mincer proposition, and show that tuition $c$ per year changes the equal-value
condition to $\ln Y(S+1) - \ln Y(S) \approx r(1 + c/Y(S))$.
\end{exercise}

\begin{exercise}[$\star\star\star$]\label{exo:b3:human-capital:11}
In the CPS comparison sample, fit log earnings on age and age squared for four
schooling groups and find the peak age of each. Why might the top-coding move the
peaks?
\end{exercise}

\begin{exercise}[$\star\star\star$]\label{exo:b3:human-capital:12}
Explain, with the \omterm{def:b3:instrumental-variables:late}{local average treatment effect}, how IV estimates can exceed OLS
estimates even when \omterm{def:b3:human-capital:ability-bias}{ability bias} is positive.
\end{exercise}

\section{Problem: The Last Year of School in Brenland}\label{sec:b3:human-capital:problem}

\begin{problem}[{Weekend problem --- the last year of school in Brenland}]\label{pb:b3:human-capital:school}
Brenland's survey (\code{load("households")}) records each head's years of schooling and
gross income; use the heads who work. A reform raises the minimum schooling from 10 to
11 years for later cohorts, and everyone who would have left earlier stays to 11; the
reformed cohorts earn 2.9 log points more. Income tax is proportional; a quarter of
the measured return is a signal; the public cost of a year of school is 30~\% of the
earnings a pupil gives up. % figdata: brenland_stats

\medskip\noindent\textbf{Part I --- the return by least squares.}
\begin{enumerate}
\item How many heads work, and what is their mean schooling?
\item Regress log income on schooling. Report the slope.
\item Add age and age squared. Report the slope and its standard error.
\item Convert the slope into a percentage premium.
\item In which direction would omitted \omterm{def:b3:human-capital:ability-bias}{ability bias} it?
\item Why does the regression use income rather than hourly earnings, and what does
that change?
\end{enumerate}

\medskip\noindent\textbf{Part II --- the reform as an instrument.}
\begin{enumerate}[resume]
\item What share of workers have 10 years of schooling or fewer?
\item What is their mean schooling, and how many years does the reform add to each?
\item Compute the \omterm{def:b3:instrumental-variables:tsls}{first stage}: the rise in mean schooling.
\item Compute the Wald estimate of the return.
\item Compare it with the least-squares estimate.
\item State the \omterm{def:b3:instrumental-variables:instrument}{exclusion restriction} and one way it could fail.
\end{enumerate}

\medskip\noindent\textbf{Part III --- whose return?}
\begin{enumerate}[resume]
\item Who are the \omterm{def:b3:instrumental-variables:compliance}{compliers}?
\item Why might their return differ from the average?
\item What would a return above OLS suggest about them?
\item Why can the reform say nothing about the return to a university degree?
\item Name a design that could.
\end{enumerate}

\medskip\noindent\textbf{Part IV --- private and social returns.}
\begin{enumerate}[resume]
\item With proportional tax and no tuition, show that the private return equals the
Wald estimate.
\item Remove the signal share: what return reflects higher output?
\item Add the public cost to the cost of the year. Compute the social return.
\item Why do taxes not enter the social return?
\item Would a higher signal share make the reform less worthwhile?
\item What else could make the social return higher than the private one?
\item Compare the social return with a 3~\% discount rate.
\item State the result: the private and the social return to the extra year, in percent.
\end{enumerate}
\end{problem}
